\section{Introduction}

Modern processors often encounter delays due to dependencies between
instructions. Memory loads are particularly important because their
completion depends on address generation, translation, cache access,
and memory-ordering constraints. Value prediction attempts to reduce
these delays by predicting an instruction's result before it becomes
available, allowing dependent operations to proceed earlier.
Incorrect predictions must be detected and handled without
compromising correctness.

This project is motivated by
``Architecting Value Prediction around In-Order Execution''
\cite{ravenel2025architecting}, which explores value prediction in
in-order processors through Last Value Prediction (LVP), ASTRA-based
load-base prediction, and V-ASTRA. V-ASTRA combines a predicted load
base with a dedicated value cache to make load values available
earlier. The paper also introduces META, which coordinates LVP and
V-ASTRA, alongside mechanisms for address compression, store
tracking, scoreboard replay, and a speculative L0 cache.

The objective is to implement and evaluate related mechanisms in
gem5 using MinorCPU. Since MinorCPU differs from the processor in
the original study, the implementation requires adaptations to its
execution, memory-access, dependency-handling, and recovery paths.

The implementation includes an LVP predictor and an ASTRA-based
load-base predictor. V-ASTRA combines the predicted base with the
instruction displacement and uses a value cache to obtain the
corresponding load value. A guarded bypass allows eligible loads to
avoid a normal memory request when address validation, cache
availability, and memory-ordering checks succeed. The project also
investigates two dependency-handling designs: Design~A releases
scoreboard dependencies after a successful bypass, while Design~B
forwards a validated value to an eligible dependent load. META is
adapted to coordinate LVP and V-ASTRA by selecting a single
prediction source for each eligible consumer and tracking that
source for recovery. An additional experimental option uses
LVP-derived provisional base information for ASTRA training and
V-ASTRA candidate generation.

The implementation is evaluated using directed RISC-V regression
tests and deterministic microbenchmarks. Validation covers predictor
behavior, value-cache correctness, memory-ordering conflicts,
dependency handling, and recovery. Performance experiments compare
no prediction (NP), LVP, V-ASTRA, and META across workloads with different address,
value, and dependency patterns, including the experimental
provisional-base modes.

The remainder of this report presents the background, implementation,
validation methodology and results, and performance evaluation. It
then discusses the limitations and differences from the original
paper and directions for future work.
